\chapter*{Úvod}
\phantomsection
\addcontentsline{toc}{chapter}{Úvod}

\ac{PMP} či \ac{DAP} jsou přenosná zařízení určená pro uložení a přehrávání médií, jako je audio, video apod. Jedná se o zařízení kapesní velikosti, napájeného akumulátorem. 
%Obrovská většina zařízení dovoluje nahrávat data do zařízení pouze prostřednictvím buďto vlastního softwaru nebo pomocí rozhraní pro připojení do osobního počítače.
Rovněž zvukové karty jsou dnes povětšině prodávány samostatně jako drahá zařízení, ačkoliv jejich hlavním úkolem je převádět digitální signál na analogový.

Cílem semestrální práce je navrhnout a realizovat rozměrem malé, kapesní zařízení, jenž bude pomocí procesoru STM32U575VIT6Q schopno z SD karty přehrávat audio soubory uložené na microSD kartě, a zároveň bude po připojení jako externí zařízení k počítači schopno reprodukovat a směřovat audio signál z počítače do sluchátek uživatele.



% Úvod studentské práce, např\,\dots

% Nečíslovaná kapitola Úvod obsahuje \uv{seznámení} čtenáře s~problematikou práce.
% Typicky se zde uvádí:
% (a) do jaké tematické oblasti práce spadá, (b) co jsou hlavní cíle celé práce a (c) jakým způsobem jich bylo dosaženo.
% Úvod zpravidla nepřesahuje jednu stranu.
% Poslední odstavec Úvodu standardně představuje základní strukturu celého dokumentu.

% Tato práce se věnuje oblasti \acs{DSP} (\acl{DSP}), zejména jevům, které nastanou při nedodržení Nyquistovy podmínky pro \ac{symfvz}.%
% \footnote{Tato věta je pouze ukázkou použití příkazů pro sazbu zkratek.}

% Šablona je nastavena na \emph{dvoustranný tisk}.
% Nebuďte překvapeni, že ve vzniklém PDF jsou volné stránky.
% Je to proto, aby důležité stránky jako např.\ začátky kapitol začínaly po vytisknutí a svázání vždy na pravé straně.
% %
% Pokud máte nějaký závažný důvod sázet (a~zejména tisknout) jednostranně, nezapomeňte si přepnout volbu \texttt{twoside} na \texttt{oneside}!